\chapter{Expressions, identités et formules}
\section*{Exercice 5.1 -- Développer}
$(2x-3)^2=4x^2-12x+9$ ; $(x+4)(x-4)=x^2-16$. Dans le premier calcul, le terme croisé est $-2\times2x\times3$.

\section*{Exercice 5.2 -- Factoriser}
$6x^2-15x=3x(2x-5)$ ; la seconde expression vaut $(x+1)(2x-3+4)=(x+1)(2x+1)$. On met d'abord le facteur commun en évidence.

\section*{Exercice 5.3 -- Domaine conservé}
Le domaine est $\R\setminus\{-3\}$. Pour ces valeurs, la factorisation $(x-3)(x+3)$ donne $x-3$. La simplification ne supprime pas l'interdiction initiale : on ne peut pas évaluer le quotient en $-3$.

\section*{Exercice 5.4 -- Isoler}
$h=V/(\pi r^2)$ puisque le dénominateur est non nul. Puis $r^2=V/(\pi h)$ ; la condition de longueur sélectionne $r=\sqrt{V/(\pi h)}$. Sans $r\geq0$, il faudrait considérer les deux signes lorsque $V>0$.

\section*{Exercice 5.5 -- Choisir la forme}
$x^2-8x+20=(x-4)^2+4$. Son minimum vaut $4$, atteint en $x=4$. L'égalité à $8$ donne $(x-4)^2=4$, donc $x=2$ ou $x=6$.

\section*{Exercice 5.6 -- Moyennes}
$(a-b)^2\geq0$ donne $(a+b)^2\geq4ab$. Comme $a+b=12$, on obtient $144\geq4ab$, donc $ab\leq36$. L'égalité impose $a=b=6$.
